\section{Introduction}

Large language models (LLMs) have demonstrated strong performance in generating fluent and contextually relevant text for a wide range of natural language processing tasks. However, their use in knowledge-intensive domains remains limited by concerns related to factual reliability, evidence grounding, and source attribution. These concerns are especially important in biomedical informatics, clinical natural language processing, healthcare document analysis, and scientific question answering, where generated responses must be supported by verifiable evidence rather than only by the internal knowledge of the model.

Retrieval-augmented generation (RAG) has emerged as an effective approach for improving the factual grounding of LLM-generated responses. Instead of relying only on information stored in model parameters, RAG systems retrieve relevant external evidence and provide that evidence as context during answer generation. Lewis et al. introduced RAG as a framework that combines parametric knowledge from a pretrained language model with non-parametric knowledge retrieved from external documents \cite{lewis2020rag}. This design allows generated responses to be guided by source documents, making RAG particularly useful for applications that require domain-specific, document-grounded, or frequently updated knowledge.

Although RAG improves the ability of LLMs to use external evidence, the quality of the final generated response depends strongly on the quality of the retrieved context. If the retrieved chunks are incomplete, weakly relevant, or poorly ranked, the generation model may produce responses that are only partially supported by the source documents. This limitation is important because semantically similar passages are not always equally useful for answering a specific query. Therefore, an effective RAG system requires not only document retrieval but also evidence prioritization.

Hybrid retrieval provides one way to improve retrieval quality by combining lexical and semantic search signals. Lexical retrieval can capture exact keyword matches, while semantic retrieval can identify conceptually related passages even when the wording differs from the query. OpenSearch supports hybrid search by combining keyword-based and vector-based retrieval signals to improve retrieval relevance \cite{opensearchhybrid}. In a RAG pipeline, this combination can increase the likelihood that the retrieved evidence includes both exact terminology and semantically relevant content.

Reranking can further improve the selection of evidence after the initial retrieval stage. A reranker receives the query and a set of candidate chunks, then reorders those chunks based on their relevance to the query. Cohere reranking models are designed to sort candidate text passages according to semantic relevance \cite{coherererank}. In the context of RAG, reranking is useful because the initially retrieved chunks may contain relevant information but may not be ordered in the most useful way for answer generation. By prioritizing the most relevant evidence before generation, reranking can help reduce unsupported or weakly grounded responses.

Another important challenge in RAG systems is evaluation. Many RAG evaluations focus on the answer as a whole, but response-level evaluation may hide unsupported statements within an otherwise acceptable answer. A generated response may include several supported statements and one unsupported factual claim. Therefore, claim-level grounding evaluation provides a more detailed and conservative assessment of response reliability. In this approach, each generated answer is decomposed into individual factual claims, and each claim is evaluated against the retrieved evidence to determine whether it is directly supported.

In this study, we propose a hybrid retrieval and reranking framework for citation-aware retrieval-augmented generation. The framework uses Amazon Bedrock Knowledge Bases to support document ingestion, parsing, chunking, embedding generation, and evidence retrieval. Amazon Bedrock Knowledge Bases allow foundation models to retrieve relevant information from user-provided data sources and use that information during response generation \cite{awsbedrockkb}. In the proposed workflow, source PDF documents are stored in Amazon S3, embedded using Amazon Titan Text Embeddings V2, and indexed using Amazon OpenSearch Serverless. The system then retrieves candidate evidence chunks, applies reranking to improve evidence prioritization, and generates answers using the highest-ranked evidence.

The proposed framework also separates answer generation from grounding evaluation. The answer-generation model produces responses using only the retrieved and reranked evidence, while a separate judge model evaluates whether each generated claim is supported by the retrieved evidence. This separation is intended to reduce self-evaluation bias and provide a more transparent assessment of grounding quality. A claim is considered supported only when the retrieved evidence directly justifies the full meaning of the claim. If the evidence is incomplete, vague, indirect, or only partially relevant, the claim is treated as unsupported.

The main contributions of this study are fourfold. First, this work presents an end-to-end RAG pipeline that integrates Amazon Bedrock Knowledge Bases, Amazon S3, Titan Text Embeddings V2, OpenSearch Serverless, hybrid retrieval, reranking, and LLM-based answer generation. Second, it introduces a claim-level grounding evaluation workflow to assess whether individual generated claims are supported by retrieved evidence. Third, it separates the answer-generation and grounding-evaluation stages by using different models for generation and judging. Fourth, it reports retrieval, reranking, query-level, and claim-level grounding results to evaluate the reliability of the proposed framework. Overall, this study aims to support the development of more reliable, transparent, and evidence-grounded RAG systems for biomedical and healthcare-related information retrieval tasks.